\section{Background}

\citet{stamatatos2009survey} provides a comprehensive survey of modern stylometric approaches that use lexical, syntactic, and structural features to identify text authorship. The survey examines both traditional statistical methods and machine learning approaches for capturing distinctive writing patterns. Although these methods achieve strong performance on controlled datasets, the robustness of stylometric features under systematic text transformations remains an open question.

\citet{koppel2011authorship} investigate authorship attribution in real-world scenarios where authors may deliberately obfuscate their writing style. They demonstrate that text transformations such as paraphrasing, translation round-trips, and stylistic rewriting significantly degrade the performance of standard authorship attribution systems. While their work establishes that obfuscation is possible, it does not systematically quantify the rate of stylometric decay across iterative transformations.

\citet{papineni2002bleu} introduce the BLEU metric, which provides an automatic method for evaluating machine translation quality through n-gram precision matching between generated and reference texts. This metric efficiently captures lexical overlap and has become the standard baseline for paraphrase evaluation in NLP. However, BLEU measures surface-level similarity rather than semantic equivalence, making it insufficient for capturing deeper authorial identity preservation.

\citet{lin2004rouge} presents ROUGE, a recall-oriented metric package originally designed for automatic summarization evaluation that measures n-gram overlap between system outputs and reference texts. ROUGE complements precision-focused metrics like BLEU by emphasizing content preservation. However, like BLEU, it remains a lexical metric and cannot distinguish between texts that share meaning but use entirely different words.

\citet{zhang2020bertscore} propose BERTScore, which evaluates text similarity using contextual BERT embeddings rather than surface n-gram overlap. BERTScore captures semantic equivalence more effectively than BLEU or ROUGE by comparing texts in embedding space. While this enables measurement of meaning preservation through paraphrasing, it was not designed to measure stylometric preservation, leaving a gap our work addresses by combining both lexical and semantic metrics to track the divergence between style decay and content preservation.